% Draft snippets for distribution-aware generator complexity
% Goal: make the "exponentially rare outputs are useless" objection explicit.
\documentclass[11pt]{article}
\usepackage[a4paper,margin=1in]{geometry}
\usepackage{amsmath,amssymb,amsthm,mathtools}
\usepackage[T1]{fontenc}
\usepackage{lmodern}
\usepackage{hyperref}

\newcommand{\Lang}[1]{\mathcal{L}(#1)}
\newcommand{\SPG}{\text{SPG}}
\newcommand{\PG}{\text{PG}}
\newcommand{\NP}{\text{NP}}
\newcommand{\SLG}{\text{SLG}}
\newcommand{\EPG}{\text{EPG}}
\newcommand{\SPGvis}[2]{\mathrm{SPG}_{\mathrm{vis}}^{#1,#2}}

\newtheorem{definition}{Definition}[section]
\newtheorem{proposition}[definition]{Proposition}
\newtheorem{corollary}[definition]{Corollary}
\theoremstyle{remark}
\newtheorem{remark}[definition]{Remark}
\newtheorem{conjecture}[definition]{Conjecture}
\theoremstyle{definition}
\newtheorem{example}[definition]{Example}

\title{Visibility-Aware Generator Complexity (Draft)}
\author{}
\date{}

\begin{document}
\maketitle
\tableofcontents
\bigskip

%% ============================================================
%% Core Framework: Target-Set Visible Mass
%% ============================================================

Throughout this section, $G$ denotes a valid generator in the sense of the
LICS paper: a halting deterministic Turing machine reading a fair-coin
bitstream.  The induced output distribution is
$D_G(x) = \mu(\{bs \in \{0,1\}^\omega : G(bs) = x\})$.
For a set $S \subseteq \Lang{G}$, write $D_G(S) = \sum_{x \in S} D_G(x)$.

\begin{definition}[Target set]\label{def:target-set}
A \emph{target set} for generator~$G$ is any set
$S \subseteq \Lang{G}$ representing outputs of interest to the tester
(e.g.\ inputs that trigger a particular bug, exercise a specific code path,
or satisfy a coverage criterion).
\end{definition}

\begin{definition}[Reachability of a target set]\label{def:reachable}
For a budget function $T : \mathbb{N} \to \mathbb{N}$, a target set~$S$ is
\emph{$T$-reachable} by~$G$ if $D_G(S) \ge 1/T(n)$, where $n$ is a size
parameter.  It is \emph{$T$-unreachable} otherwise.
\end{definition}

\begin{definition}[Visible mass of a target family]\label{def:visible-mass}
Let $\mathcal{S} = \{S_1, \dots, S_k\}$ be a finite family of target sets
(not necessarily disjoint).  The $T$-visible mass of~$G$ w.r.t.\
$\mathcal{S}$ is
\[
  \mathsf{Vis}_T(G, \mathcal{S})
  = \frac{|\{j : D_G(S_j) \ge 1/T(n)\}|}{k},
\]
the fraction of target sets that are $T$-reachable.
\end{definition}

\begin{remark}[Why target sets, not individual outputs?]\label{rem:why-target}
For structured languages with $2^{\Omega(n)}$ valid outputs of size~$n$,
every generator making $\Theta(n)$ random selections assigns each
individual output probability at most $p^{\Theta(n)} = 2^{-\Omega(n)}$
(where $p < 1$ is the maximum single-step probability).
No individual output is $T$-reachable under polynomial~$T$.

However, a \emph{set} of outputs can have polynomial or even constant
probability: the set of all inputs that trigger a particular compiler
bug aggregates exponentially many individual outputs whose probabilities
sum to a non-negligible value.
The target-set framework measures whether the generator places enough
mass on each \emph{category of interest}, not on each individual output.
\end{remark}

\begin{example}[Bug reachability]\label{ex:bug-reach}
Consider \emph{differential testing} of two compilers $C_1, C_2$
(e.g.\ GCC vs.\ Clang).  A program generator~$G$ (e.g.\ Csmith) produces
random source programs, which are then compiled and executed by both
compilers.  Here $C_i(x)$ denotes the observable output of compiler~$i$
on source program~$x$.  Define the target set
\[
  S_{\mathrm{bug}} = \{x \in \Lang{G} : C_1(x) \ne C_2(x)\}.
\]
Note that $G$ and $C_1, C_2$ play distinct roles: $G$ is the test-input
generator whose output distribution $D_G$ we analyse, while $C_1, C_2$ are
the systems under test that \emph{define} what counts as a bug.
Even if no single bug-triggering input has detectable probability under
$D_G$, the aggregate $D_G(S_{\mathrm{bug}})$ can be non-negligible because
$S_{\mathrm{bug}}$ collects many inputs.
The target set is $T$-reachable iff a disagreement can be found in
$O(T(n))$ samples from~$G$.
\end{example}

\begin{example}[Branch coverage targets]\label{ex:branch-reach}
Let $P$ be a program under test with branches $b_1, \dots, b_k$.
Define $S_j = \{x \in \Lang{G} : x \text{ exercises branch } b_j\}$.
The target family $\mathcal{S} = \{S_1, \dots, S_k\}$ captures the
coverage objectives.  $\mathsf{Vis}_T(G, \mathcal{S})$ is the fraction
of branches reachable in $T(n)$ samples.
RustSmith achieves $\mathsf{Vis}_T \approx 0.52$ on the MIR optimizer
(51.59\% line coverage), while the official test suite achieves
$\mathsf{Vis}_T \approx 0.82$ (82.12\%).
\end{example}

\begin{example}[Optimization-pass coverage]\label{ex:opt-pass}
Csmith generates C programs for differential testing of GCC/LLVM.
Define $S_j$ = inputs that exercise optimization pass~$j$.
Csmith's distribution concentrates mass on programs that trigger
constant folding, dead-code elimination, and simple inlining
($D_G(S_j)$ large for these~$j$), but assigns negligible mass to
programs triggering loop vectorization or auto-parallelization
($D_G(S_j) \approx 0$ for these~$j$).
This explains why Csmith finds many optimization bugs in simple passes
but misses bugs in complex passes entirely.
\end{example}

\begin{proposition}[Observability]\label{prop:observability}
Let $X_1,\dots,X_m$ be i.i.d.\ samples from $D_G$ and let $S$ be a
target set.  Then
\[
  \Pr[\exists i \le m.\; X_i \in S]
  = 1 - (1-D_G(S))^m.
\]
If $D_G(S) \ge 1/T(n)$, then $m = O(T(n) \log(1/\delta))$ samples
suffice to hit $S$ with probability $\ge 1-\delta$.
\end{proposition}

\begin{proof}
$\Pr[\forall i.\; X_i \notin S] = (1 - D_G(S))^m$.
If $D_G(S) \ge 1/T(n)$, set $m = T(n) \ln(1/\delta)$; then
$(1-1/T(n))^{T(n)\ln(1/\delta)} \le e^{-\ln(1/\delta)} = \delta$.
\end{proof}

\begin{corollary}[Coverage discovery curve]\label{cor:discovery}
Given target family $\mathcal{S} = \{S_1, \dots, S_k\}$, the expected
number of targets hit after $m$ samples is
\[
  \mathsf{Cov}(m) = \sum_{j=1}^{k} \bigl(1 - (1 - D_G(S_j))^m\bigr).
\]
This is precisely the species-accumulation curve from ecology, where each
target set is a ``species.''
\end{corollary}

\begin{remark}[Connection to Good-Turing estimation]\label{rem:good-turing}
In practice, $D_G(S_j)$ is unknown.  However, the \emph{discovery curve}
$\mathsf{Cov}(m)$ is observable: we can count how many distinct coverage
targets have been hit after $m$ samples.  The Good-Turing estimator
provides a nonparametric estimate of the \emph{missing mass}---the total
probability on targets not yet observed:
\[
  \widehat{M}_m = \frac{f_1(m)}{m},
\]
where $f_1(m)$ is the number of targets seen exactly once in $m$ samples
(``singletons'').  This is an unbiased estimator of the true missing mass
$\sum_{j : S_j \text{ not yet hit}} D_G(S_j)$.

When $\widehat{M}_m \to 0$, the generator has exhausted its reachable
targets: further sampling is unlikely to discover new coverage.
This gives a \emph{theory-grounded stopping criterion} for fuzzing
campaigns.
\end{remark}

\subsection{Empirical Validation}\label{sec:empirical}

We validate the Good-Turing missing mass framework experimentally using
AFL++~4.00c on a self-contained JSON parser ($\sim$190 instrumented edges).

\paragraph{Setup.}
Four campaigns, each 300 seconds:
(1)~\texttt{explore} schedule (default);
(2)~\texttt{fast} schedule;
(3)~\texttt{rare} schedule (prioritises rare edges);
(4)~\texttt{gt\_guided}: switch schedule every 100\,s
(\texttt{explore}$\to$\texttt{rare}$\to$\texttt{fast}).
After each campaign, we run \texttt{afl-showmap} on every queue entry
(in discovery order) to obtain per-input edge sets, then compute the
Good-Turing missing mass $\widehat{M}_m = f_1/m$ and Chao1 richness
estimator at each step.

\paragraph{Results.}
\begin{center}
\begin{tabular}{lcccccc}
\hline
Campaign & Corpus & Edges & 50\% at & 90\% at & Sat.\ at & $\widehat{M}_{\mathrm{final}}$ \\
\hline
\texttt{explore}    & 414 & 188 & \#83  & \#301 & \#392 & 0.073 \\
\texttt{fast}       & 418 & 187 & \#90  & \#289 & \#410 & 0.084 \\
\texttt{rare}       & 405 & 186 & \#72  & \#255 & \#399 & 0.084 \\
\texttt{gt\_guided} & 408 & 186 & \#28  & \#207 & \#372 & 0.071 \\
\hline
\end{tabular}
\end{center}

Columns ``50\% at'' and ``90\% at'' report the queue index at which 50\%
and 90\% of final edge coverage was reached; ``Sat.'' is the last queue
entry that discovered a new edge.

\paragraph{Observations.}
\begin{enumerate}
  \item \textbf{GT-guided switching reaches 90\% coverage 1.5$\times$
    faster} than the best single-schedule baseline: input~\#207 vs.\
    \#255 (\texttt{rare}).  This confirms that switching strategy at
    saturation accelerates discovery.

  \item \textbf{The discovery curve matches the theoretical model.}
    The expected coverage $\mathsf{Cov}(m) = \sum_j (1-(1-D_G(S_j))^m)$
    predicts concave-then-flat growth, which is exactly what we observe:
    rapid initial discovery followed by a long tail.

  \item \textbf{The Good-Turing missing mass is a useful real-time
    diagnostic.}  At input~\#100, $\widehat{M}$ ranges from 0.30 to 0.41
    across campaigns (``much remains to discover'').  By input~\#300, it
    drops below 0.15 (``diminishing returns'').  A fuzzer can use
    $\widehat{M} < \varepsilon$ as a principled trigger for strategy
    switching.

  \item \textbf{Chao1 overestimates total reachable edges by 24--52\%},
    consistent with the known upward bias of Chao1 when sampling is
    incomplete.  Nevertheless, the qualitative signal (``there is more to
    find'' vs.\ ``we're nearly done'') is reliable.
\end{enumerate}

\begin{proposition}[Visible support sanity]
If $\mathsf{Vis}_T(G) > 0$, then $1 \le \mathsf{VSupp}_T(G) \le |\mathsf{Heavy}_T(G)|$. Equality on the right holds iff $D_{G,T}$ is uniform on $\mathsf{Heavy}_T(G)$.
\end{proposition}

\begin{proposition}[Visible entropy bounds]
If $\mathsf{Vis}_T(G) > 0$, then
\[
\mathsf{H}^{\min}_T(G) \le \log \mathsf{VSupp}_T(G) \le \log |\mathsf{Heavy}_T(G)|.
\]
Moreover, $1 \le \mathsf{ESupp}_T(G) \le \mathsf{VSupp}_T(G)$, with equality throughout when the visible mass is uniform over the heavy outputs.
\end{proposition}

\begin{proposition}[Visible collision probability]
If $X_1$ and $X_2$ are independent samples from $D_G$ conditioned on $\mathsf{Heavy}_T(G)$, then
\[
\Pr[X_1 = X_2] = \sum_{x \in \mathsf{Heavy}_T(G)} D_{G,T}(x)^2 = \frac{1}{\mathsf{VSupp}_T(G)}.
\]
In particular, the visible support size is the inverse collision probability of the visible outputs.
\end{proposition}

\begin{corollary}[Expected visible duplicates]
For $m$ independent samples from $D_G$ conditioned on $\mathsf{Heavy}_T(G)$, the expected number of colliding pairs is
\[
\binom{m}{2} / \mathsf{VSupp}_T(G).
\]
So large visible support size means repeated visible outputs are rare.
\end{corollary}

\begin{proposition}[Invisible outputs under finite budgets]
If an output $x$ has probability at most $\varepsilon$, then the probability that $x$ appears at least once in $m$ independent samples is at most $m\varepsilon$.
In particular, if $\varepsilon = 2^{-\Omega(n)}$ and $m = \mathrm{poly}(n)$, then $x$ is practically invisible.
\end{proposition}

\begin{remark}
This is a refinement of the support-based view rather than a replacement for it: a generator may be support-correct but still have vanishing visible mass. For fuzzing, that means validity alone is not enough; the generator must also place enough mass on outputs that can be reached under realistic sampling budgets. Equivalently, the distribution should not have a tiny effective support size relative to the actual support.
\end{remark}

% Possible follow-up theorem direction:
% Characterize which languages admit polynomial-time visible-mass generators.
% Compare support-wise GEN with visibility-aware GEN^{vis}_{T,\rho}.
% Relate visibility to coverage / throughput metrics in fuzzing practice.

%% ============================================================
%% Phase 1 & 3 exposition: PBT basics, generator combinators,
%% SPG connection, RustSmith mapping
%% ============================================================

\section{Background: Property-Based Testing and Generators}\label{sec:pbt-background}

\subsection{Property-Based Testing vs Unit Testing}

Unit testing checks a small set of hand-picked input--output examples:
\begin{verbatim}
  assert add(2, 3) == 5
  assert add(0, 0) == 0
\end{verbatim}
Property-based testing (PBT) instead states a \emph{property}---a universal
quantified statement that must hold for all inputs---and uses a \emph{generator}
to produce many random inputs automatically:
\begin{verbatim}
  @given(integers(), integers())
  def test_add_commutative(x, y):
      assert add(x, y) == add(y, x)
\end{verbatim}
When a counterexample is found, PBT frameworks \emph{shrink} it to a minimal
failing input.

\subsection{Generator Combinators}

In frameworks like QuickCheck~\cite{claessen2000quickcheck} and
Hypothesis~\cite{MacIver2019}, generators are built compositionally from a small
set of combinators:

\begin{description}
  \item[Atomic generators] \texttt{choose(lo, hi)}: uniform random integer in
    $[\texttt{lo}, \texttt{hi}]$; \texttt{pure(a)}: always return \texttt{a}.
  \item[Mapping] \texttt{fmap(f, g)}: apply function $f$ to the output of
    generator $g$.
  \item[Choice] \texttt{oneof([g1, ..., gk])}: pick uniformly among $k$
    generators. \texttt{frequency([(w1, g1), ..., (wk, gk)])}: weighted choice.
  \item[Collection] \texttt{listOf(g)}: generate a list of random length whose
    elements come from $g$.
  \item[Recursion] For recursive data types (e.g.\ ASTs), one uses
    \texttt{frequency} with a leaf case and a recursive case, where the
    recursive case calls the generator again for subtrees.
\end{description}

Each combinator has distributional implications.  For instance, if a branch in
\texttt{frequency} has weight $w$ out of total $W$, its probability is $w/W$;
the expected recursion depth to reach a leaf is $O(W/w)$.  If $w/W = 2^{-\Omega(n)}$, then deep structures appear with exponentially small
probability---precisely the ``exponentially rare outputs'' problem.

\subsection{SPG: Size-Dependent Polynomial Generators}\label{sec:spg-recap}

The LICS paper defines SPG (Definition~4.6) as follows.  A \emph{sized
partition}\label{def:size-partition} of a language $L$ is a partition $(L_i)_{i \in C}$ indexed by a set
$C \subseteq \{0,1\}^*$ with a size function $l : C \to \mathbb{N}$ and
polynomial-time inverse $l^{-1}$.  A \emph{size-dependent generator}
$G = (G_i)_{i \in C}$ satisfies $\mathcal{L}(G_i) = L_i$ for each $i \in C$.

\[
\SPG{} = \{\mathcal{L}(G) : G \text{ is an } l\text{-sized polynomial-time
generator for some partition}\}.
\]

Key results: $\SPG{} \subseteq \PG{}$, $\SPG{} \subseteq \NP{}$,
$\mathsf{CNF\text{-}SAT} \in \SPG{}$, and under collision-resistant hash
assumptions $\SPG{} \subsetneq \NP{}$.

\subsection{Mapping RustSmith to SPG}\label{sec:rustsmith-spg}

RustSmith~\cite{sharma2023rustsmith} builds random, well-typed Rust programs
constructively (no rejection sampling).  Key observations:

\begin{enumerate}
  \item \textbf{SPG membership:} RustSmith's generator is polynomial-time in
    output size (symbol table lookups, context-aware AST construction).  Its
    language $L_{\text{RustSmith}} \subseteq L_{\text{well-typed}}$ is in SPG.
  \item \textbf{Type~1 errors:} RustSmith deliberately under-generates
    (conservative borrowing, fixed lifetimes).  This is
    $L_{\text{RustSmith}} \subsetneq L_{\text{well-typed}}$.
  \item \textbf{Zero Type~2 errors:} Every generated program compiles.  No
    over-generation.
  \item \textbf{Distribution matters:} Coverage experiments show some
    optimization passes get 100\% coverage, others 0\%.  The generator's output
    distribution controls practical utility, not just support correctness.
\end{enumerate}

This motivates the visible-mass framework: SPG ensures support coverage, but
visible mass captures whether the ``useful'' outputs are actually reachable under
finite sampling budgets.

%% ============================================================
%% Phase 1.5: Engineering challenges
%% ============================================================

\subsection{Key Engineering Challenges}\label{sec:eng-challenges}

Three intertwined difficulties arise when building a practical generator for a
structured language:

\begin{description}
  \item[Type conformance] The generator must produce inputs satisfying the
    language's typing rules.  For simple type systems (e.g.\ C without
    generics), constructive generation is polynomial-time.
  \item[Semantic validity] The generator must avoid undefined behavior.
    RustSmith uses wrappers to guard dangerous operations.  Exact generation of
    $L_{\text{well-defined}}$ is in general undecidable (a Rice-theorem-style
    argument).
  \item[Distribution control] The generator must steer probability mass toward
    useful structures.  Rejection sampling is too expensive for structured
    inputs; the generator must be \emph{constructive}.
\end{description}

RustSmith exemplifies the pragmatic choice: accept Type~1 errors
(under-generation) to achieve constructive generation in SPG, while maintaining
zero Type~2 errors.  The visible-mass framework quantifies the cost of this
choice: $\mathsf{Vis}_T(G)$ and $\mathsf{VSupp}_T(G)$ measure how much of the
target language is practically reachable.

%% ============================================================
%% Phase 3.1--3.3: Where does well-typed Rust sit?
%% ============================================================

\subsection{Well-Typed Rust in the Complexity Hierarchy}\label{sec:rust-hierarchy}

Let $L_{\text{Rust}}$ = syntactically valid Rust, $L_{\text{well-typed}}$ =
programs passing the type checker, $L_{\text{well-defined}}$ = programs with
well-defined runtime behavior.

\[
L_{\text{well-defined}} \subseteq L_{\text{well-typed}} \subseteq L_{\text{Rust}}.
\]

\begin{description}
  \item[$L_{\text{Rust}} \in \SPG$:] Syntactic generation from a grammar is
    polynomial-time.
  \item[$L_{\text{well-typed}} \in \SPG$?] Likely yes for Rust without
    generics/traits (type checking is P-time).  With trait resolution (which
    encodes logic programming), complexity may be higher.
  \item[$L_{\text{well-defined}} \in \SPG$?] Unlikely in general: deciding
    well-defined runtime behavior is undecidable.
\end{description}

\subsection{Type~1 vs Type~2 Error Tradeoffs}\label{sec:type1-type2}

Define parameterised generation classes:
\[
\SPG_{\alpha,\beta} = \{ \mathcal{L}(G) : G \in \SPG{},\;
\Pr[\text{Type~1 error}] \le \alpha,\;
\Pr[\text{Type~2 error}] \le \beta \}.
\]
RustSmith's strategy corresponds to $\beta = 0$ (zero over-generation) with
$\alpha$ determined by its conservative borrow-checking.  In differential
testing, Type~2 errors (over-generation) produce false-positive bug reports and
waste human review; Type~1 errors merely reduce coverage.

%% ============================================================
%% Phase 3.4: Why distribution matters
%% ============================================================

\subsection{Why Distribution Matters: The Visible-Mass Bridge}\label{sec:distribution-matters}

If generator $G$ has output distribution $D_G$ and a target set
$T \subseteq \mathcal{L}(G)$ of ``useful'' outputs, the expected number of
samples to first hit $T$ is $1/D_G(T)$.  If $D_G(T) = 2^{-\Omega(n)}$, then
$2^{\Omega(n)}$ samples are needed---practically impossible.

This is the gap between SPG and practice that visible mass fills:
\begin{itemize}
  \item SPG guarantees $\mathrm{supp}(D_G) = \mathcal{L}(G)$.
  \item Visible mass asks: ``Of this support, how much probability is
    concentrated on outputs that can actually be observed under budget $T$?''
  \item A generator in SPG with $\mathsf{Vis}_T(G) \to 0$ is correct in theory
    but useless in practice.
\end{itemize}

%% ============================================================
%% Phase 4: Reviewer Response Formalizations
%% ============================================================

\section{Distribution-Aware Generator Classes}\label{sec:distribution-aware}

\subsection{Visibility-Restricted Generation Classes}\label{sec:vis-classes}

We now define distribution-aware complexity classes that address the critique
that support-only generability ignores exponentially rare outputs.

\begin{definition}[Visibility-restricted class]
For a budget function $T : \mathbb{N} \to \mathbb{N}$ and a visibility
threshold $\rho > 0$, define
\[
\SPGvis{T}{\rho} = \{ \mathcal{L}(G) : G \in \SPG{},\;
\mathsf{Vis}_T(G) \ge \rho \}.
\]
A language is in $\SPGvis{T}{\rho}$ if it has a size-dependent polynomial-time
generator that concentrates at least $\rho$ fraction of its probability mass on
$T$-heavy outputs.
\end{definition}

\begin{proposition}[Continuous hierarchy]
For fixed $T$, the classes $\SPGvis{T}{\rho}$ form a decreasing family:
if $\rho_1 \le \rho_2$, then
$\SPGvis{T}{\rho_1} \supseteq \SPGvis{T}{\rho_2}$.
\end{proposition}

\begin{proposition}[SPG-vis separation from SPG]\label{prop:spg-vis-separation}
There exist languages in $\SPG{} \setminus \SPGvis{T}{1/\mathrm{poly}(n)}$.
Specifically, let $L$ be a language with $|L \cap \{0,1\}^n| = 2^n$ and consider
a generator $G$ that assigns probability $2^{-n}$ to each output.  Then
$G \in \SPG{}$ but $\mathsf{Vis}_T(G) = 0$ for $T(n) = \mathrm{poly}(n)$.
\end{proposition}

\begin{remark}
This separation justifies the reviewer concern: a language can be
generable in polynomial time (support-correct) yet practically unreachable
under any realistic sampling budget.  The $\SPGvis{T}{\rho}$ hierarchy
captures the ``practically generable'' languages.
\end{remark}

\subsection{Practical Complexity Metrics}\label{sec:practical-metrics}

We formalise three metrics used in fuzzing practice.

\begin{definition}[Rejection rate]
For a generator $G$ that may produce outputs outside the target language $L$,
define the rejection rate as
\[
\beta(G) = 1 - \Pr_{b}[G(b) \in L].
\]
\end{definition}

\begin{definition}[Effective throughput]
Let $\mathrm{time}(G)$ be the expected time per output of $G$.  The effective
throughput is
\[
\mathrm{TP}(G) = \frac{1 - \beta(G)}{\mathrm{time}(G)}.
\]
A generator with $\beta(G) = 1 - \varepsilon$ has throughput
$\varepsilon / \mathrm{time}(G)$, which is exponentially small if rejection
sampling is used for a language where valid outputs are rare.
\end{definition}

\begin{definition}[Coverage rate]
For budget $T$ and sample size $m$, the coverage rate of $G$ is
$\mathsf{Occ}_{m,T}(G) / |\mathsf{Heavy}_T(G)|$ (fraction of heavy outputs
discovered).  Under the visibility-rich condition
$(\rho, \sigma)$, Corollary~\ref{cor:coverage-lb} gives a lower bound.
\end{definition}

\subsection{Type~1 vs Type~2 Error Classes}\label{sec:type1-type2-formal}

\begin{definition}[Error-parameterised classes]
For $\alpha, \beta \in [0,1]$, define
\[
\SPG_{\alpha,\beta} = \{ \mathcal{L}(G) : G \in \SPG{},\;
\Pr[\text{Type~1 error}] \le \alpha,\;
\Pr[\text{Type~2 error}] \le \beta \}.
\]
Here a Type~1 error means a valid output $w \in L$ has
$D_G(w) = 0$ (under-generation), and a Type~2 error means an output
$w \notin L$ has $D_G(w) > 0$ (over-generation).
\end{definition}

\begin{conjecture}[Strict separation via error relaxation]
There exist languages $L$ such that $L \notin \SPG_{0,0}$ (no zero-error
polynomial-time generator exists) but $L \in \SPG_{\alpha,0}$ for some
$\alpha > 0$ (a polynomial-time generator exists that allows under-generation
but not over-generation).  The language $L_{\text{well-typed}}$ of
well-typed Rust programs is a candidate, since exact generation requires
data-flow analysis (potentially beyond SPG) while conservative approximation
is in SPG with Type~1 errors only.
\end{conjecture}

\subsection{Real-World Generator Taxonomy}\label{sec:real-world-taxonomy}

We map several well-known generators to our framework:

\begin{center}
\begin{tabular}{lcccc}
\hline
\textbf{Tool} & \textbf{Class} & \textbf{Type~1} & \textbf{Type~2} & \textbf{Distribution} \\
\hline
QuickCheck & SPG & low (filter) & low & engineered \\
Csmith     & SPG & low (wrappers) & zero & hand-tuned \\
RustSmith  & SPG & high (conservative) & zero & heavily engineered \\
AFL        & PG  & N/A & high (byte-level) & coverage-guided \\
\hline
\end{tabular}
\end{center}

QuickCheck uses filter-based validity (low Type~1, some Type~2 from incomplete
specifications).  Csmith wraps dangerous operations (low Type~1, zero Type~2).
RustSmith deliberately over-constrains borrowing rules (high Type~1, zero
Type~2).  AFL operates on raw bytes with no validity constraint (high Type~2
but reaches arbitrary inputs via mutation).

%% ============================================================
%% Phase 5: Related Work and Connections
%% ============================================================

\section{Related Work}\label{sec:related-work}

\subsection{Samplable Distributions and Average-Case Complexity}

Levin~\cite{levin1986average} introduced the theory of average-case complexity
and the notion of P-samplable distributions: a distribution is P-samplable if a
polynomial-time TM can draw samples from it.  Our generator complexity classes
are a refinement: we ask not just ``can you sample?'' but ``what is the
complexity of sampling with good coverage?''

Impagliazzo and Levin~\cite{impagliazzo1990better} showed that if a distribution
is easy to sample, then any algorithm solving NP problems under that
distribution implies a randomized algorithm for all of NP.  Watson~\cite{watson2017time}
proved strict time hierarchies for sampling: there exist distributions
samplable in time $O(f(n))$ but not in $o(f(n))$.  This directly applies: different
generators for the same language may provably differ in sampling speed.

\subsection{Enumeration Complexity}

Arenas, Jerrum, Segoufin, \& Sirangelo~\cite{arenas2019constant} showed that for
conjunctive queries, there exist constant-delay enumeration algorithms after
polynomial preprocessing.  The delay-based complexity measure is analogous to our
generator throughput: the cost per additional output.

The enumeration hierarchy---polynomial preprocessing + constant delay vs.\
polynomial delay vs.\ incremental polynomial time---maps naturally to generator
complexity classes.  Our SPG corresponds to polynomial preprocessing; the delay
corresponds to the per-sample cost.

\subsection{Random Generation and Counting}

Jerrum, Valiant, \& Vazirani~\cite{jerrum1986random} established the deep
connection between random generation and approximate counting: almost-uniform
generation and approximate counting are polynomially related.  This implies that
if counting valid programs is \#P-hard, then uniform generation is also hard.

Sanchis~\cite{sanchis1990data} defined the predecessor of our SPG notion: a
generator is efficient if its expected time per output is polynomial in output
size.  The main result---that systematic generation from algebraic
specifications can be done in polynomial time---directly supports our claim
that $L_{\text{Rust}} \in \SPG{}$ (syntactic generation is efficient).

\subsection{Boltzmann Sampling}

Duchon, Flajolet, Louchard, \& Schaeffer~\cite{duchon2004boltzmann} introduced
Boltzmann samplers: each structure of size $n$ gets weight proportional to $x^n$,
and the parameter $x$ controls the expected size.  This is a principled way to
control output distribution while maintaining efficiency.

Boltzmann sampling is size-oblivious: it generates a structure and accepts if
the size is in range.  The parameter $x$ gives continuous control over the
distribution, analogous to our budget parameter $T$.  The key difference is
that Boltzmann samplers assume a context-free grammar specification, while our
framework handles arbitrary constraint-based generation.

\subsection{Coverage-Guided Fuzzing}

B\"ohme, Pham, \& Roychoudhury~\cite{bohme2016coverage} model AFL-style fuzzing
as a Markov chain over program states.  The fuzzer's effectiveness is determined
by the stationary distribution.  Their ``directed greybox
fuzzing''~\cite{bohme2017directed} extension biases the generator toward
specific targets---the practical version of our ``heavy outputs'' concept.

Lemieux \& Sen~\cite{lemieux2018fairfuzz} (FairFuzz) directly address the
``rare output'' problem: inputs exercising rare branches are
under-represented in the fuzzer's distribution, so FairFuzz reweights to
increase their probability.  This is exactly our visible-mass argument in
practice---FairFuzz is engineering $\mathsf{Vis}_T(G)$ upward.

\subsection{Diversity Measures and Effective Support}

Hill~\cite{hill1973diversity} defined the Hill number family
$H_q = (\sum p_i^q)^{1/(1-q)}$, which unifies species richness ($q=0$),
exponential Shannon entropy ($q=1$), and inverse Simpson index ($q=2$).
Our visible diversity profile is a thresholded Hill-number profile restricted
to heavy outputs.

Good~\cite{good1953population} introduced the missing mass: the probability of
seeing a new species on the next sample.  In our framework, the missing mass is
$1 - \mathsf{Vis}_T(G)$---the total probability on invisible outputs.

%% ============================================================
%% Phase 5: New Theorems and Separation Results
%% ============================================================

\section{New Results}\label{sec:new-results}

\subsection{Visible-Mass Optimality}

\begin{proposition}[Visible mass is maximised by greedy allocation]
Fix a budget function $T$ and a language $L$ with $|L \cap \{0,1\}^n| = N$.
Among all distributions $D$ with $\mathrm{supp}(D) = L \cap \{0,1\}^n$, the
visible mass $\mathsf{Vis}_T(D)$ is maximised by placing probability mass
$1/T(n)$ on each of the $\min(N, \lfloor T(n) \rfloor)$ heaviest outputs,
with the remainder on a single additional output if $N > T(n)$.
\end{proposition}

\begin{proof}
The heavy set $\mathsf{Heavy}_T$ consists of outputs $x$ with
$D(x) \ge 1/T(n)$.  Since $\sum D(x) = 1$, at most $T(n)$ outputs can
satisfy this constraint.  The visible mass is
$\sum_{x \in \mathsf{Heavy}_T} D(x)$, which is maximised by concentrating
as much mass as possible on outputs that meet the threshold.  Placing
$1/T(n)$ on each of $k = \min(N, \lfloor T(n) \rfloor)$ outputs and the
remainder on one more output achieves the maximum, which is
$\min(1, k/T(n)) = 1$ when $k \ge T(n)$.
\end{proof}

\begin{corollary}[Budget hardness]
For a language $L$ with $|L \cap \{0,1\}^n| = 2^{\Omega(n)}$ and budget
$T(n) = \mathrm{poly}(n)$, any generator $G$ for $L$ with $\mathrm{supp}(D_G) = L$
satisfies $\mathsf{Vis}_T(G) \le \mathrm{poly}(n) / 2^{\Omega(n)} = 2^{-\Omega(n)}$.
Thus $L \notin \SPGvis{T}{1/\mathrm{poly}(n)}$.
\end{corollary}

\begin{remark}
This confirms the reviewer's intuition: for exponentially large languages
(e.g., all well-typed programs of size $n$), any generator that truly covers
the full language necessarily has vanishing visible mass.  Practical generators
\emph{must} under-generate (accept Type~1 errors) to achieve non-trivial
visible mass.
\end{remark}

\subsection{Separation and Closure Properties of $\SPGvis{T}{\rho}$}\label{sec:spg-vis-closure}

We establish structural properties of the visibility-restricted class.

\begin{proposition}[Closure under intersection with decidable languages]\label{prop:closure-intersection}
Let $L \in \SPGvis{T}{\rho}$ via generator $G$, and let $S$ be decidable in
polynomial time.  Then $L \cap S \in \SPGvis{T'}{\rho'}$ for some
$T' = \mathrm{poly}(n)$ and $\rho' > 0$, provided
$D_G(L \cap S) \ge 1/\mathrm{poly}(n)$.
\end{proposition}

\begin{proof}
Construct generator $G'$ that runs $G$ and accepts only when the output is in
$S$ (rejection sampling against $S$).  Since $S \in P$, each check takes
polynomial time.  The expected number of trials is $1/D_G(L \cap S)$, which
is polynomial by assumption.  The heavy outputs of $G'$ within $L \cap S$ are
exactly the heavy outputs of $G$ that also lie in $S$, rescaled by the
acceptance probability.  A union bound gives $\mathsf{Vis}_{T'}(G') \ge
\rho \cdot D_G(S \mid \mathsf{Heavy}_T(G))$, which is bounded below by a
polynomial fraction when $S$ captures a non-negligible portion of the heavy
set.
\end{proof}

\begin{proposition}[Non-closure under union]\label{prop:non-closure-union}
There exist languages $L_1, L_2 \in \SPGvis{T}{\rho}$ such that
$L_1 \cup L_2 \notin \SPGvis{T}{\rho/2}$ for the ``natural'' combined
generator.
\end{proposition}

\begin{proof}
Let $|L_1 \cap \{0,1\}^n| = |L_2 \cap \{0,1\}^n| = T(n)$, with $L_1 \cap L_2 = \emptyset$.
Generators $G_1, G_2$ that are uniform on $L_1, L_2$ respectively achieve
$\mathsf{Vis}_T(G_i) = 1$.  Consider the generator $G$ for $L_1 \cup L_2$ that
runs $G_1$ with probability $1/2$ and $G_2$ with probability $1/2$.  Each
output now has probability $1/(2T(n))$, so $\mathsf{Heavy}_{T}(G) = \emptyset$
and $\mathsf{Vis}_T(G) = 0$.  To recover visible mass, the budget must be
doubled to $T' = 2T$, under which $\mathsf{Vis}_{T'}(G) = 1$.  This shows
the class is not closed under union at a fixed budget---the budget must grow
with the number of components.
\end{proof}

\begin{proposition}[Strict hierarchy in $\rho$]\label{prop:strict-rho-hierarchy}
For every $0 < \rho_1 < \rho_2 \le 1$ and every polynomial budget $T$,
$\SPGvis{T}{\rho_2} \subsetneq \SPGvis{T}{\rho_1}$.
\end{proposition}

\begin{proof}
The containment $\SPGvis{T}{\rho_2} \subseteq \SPGvis{T}{\rho_1}$ is
immediate.  For strict separation, fix $T$ and consider a family of languages
$L_k$ with $|L_k \cap \{0,1\}^n| = k \cdot T(n)$ for integer $k \ge 1$.
A generator $G_k$ that is uniform on $L_k$ gives each output probability
$1/(kT(n))$.  Then $\mathsf{Heavy}_T(G_k) = L_k$ iff $k = 1$; for $k > 1$,
$\mathsf{Heavy}_T(G_k) = \emptyset$ and $\mathsf{Vis}_T(G_k) = 0$.  However,
any generator for $L_k$ with $\mathrm{supp} = L_k$ satisfies
$\mathsf{Vis}_T \le 1/k$ (since it must spread mass over $kT(n)$ elements).
Choosing $k$ between $\lceil 1/\rho_2 \rceil$ and $\lfloor 1/\rho_1 \rfloor$
witnesses $L_k \in \SPGvis{T}{\rho_1} \setminus \SPGvis{T}{\rho_2}$.
\end{proof}

\begin{proposition}[Strict hierarchy in $T$]\label{prop:strict-T-hierarchy}
For budget functions $T_1(n) < T_2(n)$ (both polynomial) and fixed
$\rho > 0$, $\SPGvis{T_1}{\rho} \subsetneq \SPGvis{T_2}{\rho}$.
\end{proposition}

\begin{proof}
Containment follows from monotonicity of heavy sets:
$\mathsf{Heavy}_{T_1}(G) \subseteq \mathsf{Heavy}_{T_2}(G)$.  For strict
separation, consider a language $L$ with $|L \cap \{0,1\}^n|$ between $T_1(n)$
and $T_2(n)$.  The uniform generator on $L$ gives each output probability
$1/|L_n|$, which satisfies $\ge 1/T_2(n)$ but $< 1/T_1(n)$.  So
$\mathsf{Vis}_{T_2} = 1$ but $\mathsf{Vis}_{T_1} = 0$, and hence
$L \in \SPGvis{T_2}{\rho} \setminus \SPGvis{T_1}{\rho}$.
\end{proof}

\begin{proposition}[Relationship to EPG]\label{prop:spgvis-epg}
$\SPGvis{T}{\rho} \subseteq \EPG{}$ for any polynomial $T$ and
$\rho \ge 1/\mathrm{poly}(n)$.
\end{proposition}

\begin{proof}
Let $G$ be the SPG generator witnessing $L \in \SPGvis{T}{\rho}$.  Since
$\mathsf{Vis}_T(G) \ge \rho \ge 1/\mathrm{poly}(n)$, at least a
$1/\mathrm{poly}(n)$ fraction of probability mass lands on outputs reachable
within the polynomial time bound.  Running $G$ repeatedly until a heavy output
is produced takes expected $1/\rho = \mathrm{poly}(n)$ trials, each costing
polynomial time.  The composed generator runs in expected polynomial time and
produces exactly $\mathcal{L}(G) \cap \mathsf{Heavy}_T(G)$.  Since every heavy
output is reachable within polynomial time, this gives an EPG generator for
the language restricted to heavy outputs.  For the full language, the original
SPG generator already shows $L \in \SPG{} \subseteq \PG{}$; the EPG containment
of the visibility-restricted sublanguage follows directly.
\end{proof}

\subsection{Visible Mass and Coverage Time}

\begin{proposition}[Coverage-time phase transition]
For a $(T, \rho, \sigma)$-visibility-rich generator $G$ and sample budget
$m$:
\begin{enumerate}
  \item If $m \ll \sigma/\rho$, then $\mathsf{Occ}_{m,T}(G) \approx m\rho$
    (linear discovery regime).
  \item If $m \gg (\sigma/\rho) \ln \sigma$, then
    $\mathsf{Occ}_{m,T}(G) \approx \sigma$ (saturation regime).
  \item The transition happens at $m^* \sim (\sigma/\rho) \ln \sigma$.
\end{enumerate}
\end{proposition}

\begin{proof}
From the coverage lower bound:
$\mathsf{Occ}_{m,T}(G) \ge \sigma(1 - (1-\rho/\sigma)^m)$.
For $m \ll \sigma/\rho$, the first-order Taylor expansion gives
$\sigma \cdot m\rho/\sigma = m\rho$.
For $m \gg (\sigma/\rho)\ln\sigma$, the term $(1-\rho/\sigma)^m \to 0$,
giving $\mathsf{Occ}_{m,T}(G) \to \sigma$.
The transition point is where $(1-\rho/\sigma)^m = 1/2$, i.e.,
$m = \ln 2 / (-\ln(1-\rho/\sigma)) \approx (\sigma/\rho)\ln 2$.
\end{proof}

\subsection{Connection to Enumeration Complexity}

\begin{proposition}[Enumeration-generability implies visibility]
Let $L$ have a constant-delay enumeration algorithm after polynomial
preprocessing, and let $G$ be the generator that uses the enumeration
algorithm to produce a uniform random output of each size.  Then
$G \in \SPGvis{T}{1}$ for $T(n) = \mathrm{poly}(n)$, i.e., the generator
has full visible mass.
\end{proposition}

\begin{proof}
A constant-delay enumeration algorithm produces each output in $O(1)$ time
after polynomial preprocessing.  If the generator can enumerate all outputs
of size $n$ in polynomial time, it can assign non-zero probability to every
output and choose uniformly.  Since all outputs are reachable in polynomial
time, all have probability $\ge 1/|L_n| \ge 1/\mathrm{poly}(n)$, so all are
$T$-heavy, giving $\mathsf{Vis}_T(G) = 1$.
\end{proof}

\begin{remark}
This provides a positive direction: languages with efficient enumeration
have visibility-rich generators.  The interesting case is languages where
enumeration is harder---the visible mass drops as the cost per output
increases.
\end{remark}

%% ============================================================
%% Connection to Reviewer Concerns
%% ============================================================

\section{Addressing Reviewer Concerns}\label{sec:reviewer-response}

\subsection{Distribution-Agnoticism (Reviewer~A)}

Reviewer~A's primary concern is that our original framework is
distribution-agnostic: a generator might cover a language in the support sense
while assigning exponentially small probability to most outputs, making them
practically unreachable.

Our response: the $\SPGvis{T}{\rho}$ hierarchy (Section~\ref{sec:vis-classes})
directly addresses this.  A language in $\SPGvis{T}{\rho}$ has a generator
with provably non-trivial visible mass under budget $T$.  The separation
result (Proposition~\ref{prop:spg-vis-separation}) confirms that this is a
strict refinement of SPG.

\subsection{Output-Based Complexity (Reviewer~B)}

Reviewer~B questioned measuring time by output size, noting that a generator
outputting $2^n$ bits in $(2^n)^2$ steps is ``polynomial in output size'' but
exponential in input size.

Our response: the sized-partition framework (Definition~\ref{def:size-partition})
addresses this by parameterising generators by size labels, not output
length.  The size function $l$ can encode any notion of ``relevant size''
(number of clauses, graph size, etc.), and the polynomial bound applies to
$l(i)$, not $|w|$.  This cleanly separates the size parameter from the
encoding.

\subsection{Transducer Connection (Reviewer~C)}

Reviewer~C observed that our generators are special cases of transducers.

Our response: we acknowledge the connection explicitly and note that our
contribution is not the transducer model itself but the \emph{complexity
hierarchy} it induces for PBT.  The SPG class, the visibility-restricted
refinements, and the separation results are specific to the PBT setting
and have no analogue in classical transducer theory.

\subsection{RE $=$ GEN (Reviewer~B)}

Reviewer~B found the $\mathrm{RE} = \mathrm{GEN}$ result trivial.

Our response: we agree this is a basic observation and reframe it as a
warm-up establishing the setting.  The non-trivial content is in the
\emph{efficient} generation hierarchy: $\SPG{} \subsetneq \NP{}$ (under
cryptographic assumptions), the visibility-restricted refinements, and the
connection to enumeration complexity.

\begin{definition}[Visible occupancy profile]
For an integer $m \ge 1$, define the $m$-sample visible occupancy profile of $G$ by
\[
\mathsf{Occ}_{m,T}(G) = \sum_{x \in \mathsf{Heavy}_T(G)} \bigl(1 - (1 - D_G(x))^m\bigr).
\]
\end{definition}

\begin{definition}[Visibility-rich generators]
A valid generator $G$ is $(T,\rho,\sigma)$-visibility-rich if $\mathsf{Vis}_T(G) \ge \rho$ and $\mathsf{VSupp}_T(G) \ge \sigma$.
\end{definition}

\begin{proposition}[Visible occupancy profile]
Let $X_1,\dots,X_m$ be independent samples from $D_G$. Then $\mathsf{Occ}_{m,T}(G)$ is exactly the expected number of distinct outputs in $\mathsf{Heavy}_T(G)$ that appear among $X_1,\dots,X_m$.
Moreover, $\mathsf{Occ}_{1,T}(G)=\mathsf{Vis}_T(G)$, $\mathsf{Occ}_{m,T}(G)$ is nondecreasing in $m$ and $T$, and
\[
\mathsf{Occ}_{m,T}(G) \ge m\mathsf{Vis}_T(G) - \binom{m}{2}\frac{\mathsf{Vis}_T(G)^2}{\mathsf{VSupp}_T(G)}.
\]
In particular, if $m \le \mathsf{VSupp}_T(G)$, then
\[
\mathsf{Occ}_{m,T}(G) \ge \frac{m}{2}\mathsf{Vis}_T(G).
\]
\end{proposition}

\begin{corollary}[High-diversity regime]
If $G$ is $(T,\rho,\sigma)$-visibility-rich and $m \le \sigma$, then
\[
\mathsf{Occ}_{m,T}(G) \ge m\rho/2.
\]
So in the regime where visible mass is non-negligible and visible support is large, polynomially many samples produce polynomially many distinct heavy outputs.
\end{corollary}

\begin{definition}[Visible discovery rate]
For an integer $m \ge 0$, define the $m$-sample visible discovery rate of $G$ by
\[
\mathsf{Disc}_{m,T}(G) = \sum_{x \in \mathsf{Heavy}_T(G)} D_G(x)\bigl(1 - D_G(x)\bigr)^m.
\]
\end{definition}

\begin{proposition}[Discovery curve]
Let $X_1,\dots,X_m$ be independent samples from $D_G$. Then $\mathsf{Disc}_{m,T}(G)$ is the expected number of heavy outputs in $\mathsf{Heavy}_T(G)$ that are seen for the first time at sample $m+1$. Equivalently,
\[
\mathsf{Occ}_{m+1,T}(G) - \mathsf{Occ}_{m,T}(G) = \mathsf{Disc}_{m,T}(G).
\]
Consequently,
\[
\mathsf{Occ}_{m,T}(G)=\sum_{j=0}^{m-1}\mathsf{Disc}_{j,T}(G)
\quad\text{and}\quad
\mathsf{Disc}_{0,T}(G)=\mathsf{Vis}_T(G).
\]
Moreover, $\mathsf{Disc}_{m,T}(G)$ is nonincreasing in $m$ and satisfies
\[
\mathsf{Disc}_{m,T}(G) \ge \mathsf{Vis}_T(G)\Bigl(1-\frac{\mathsf{Vis}_T(G)}{\mathsf{VSupp}_T(G)}\Bigr)^m.
\]
\end{proposition}

\begin{corollary}[Coverage lower bound]\label{cor:coverage-lb}
If $m \ge 1$, then
\[
\mathsf{Occ}_{m,T}(G) \ge \mathsf{VSupp}_T(G)\Bigl(1 - \Bigl(1 - \frac{\mathsf{Vis}_T(G)}{\mathsf{VSupp}_T(G)}\Bigr)^m\Bigr).
\]
In particular, if $G$ is $(T,\rho,\sigma)$-visibility-rich, then
\[
\mathsf{Occ}_{m,T}(G) \ge \sigma\Bigl(1 - \Bigl(1 - \frac{\rho}{\sigma}\Bigr)^m\Bigr).
\]
The bound is tight when $D_{G,T}$ is uniform on $\mathsf{Heavy}_T(G)$.
\end{corollary}

\begin{corollary}[Sample complexity for completeness]
If $0 < \varepsilon < 1$ and
\[
 m \ge \frac{\mathsf{VSupp}_T(G)}{\mathsf{Vis}_T(G)}\ln\frac{1}{\varepsilon},
\]
then
\[
\mathsf{Occ}_{m,T}(G) \ge (1-\varepsilon)\mathsf{VSupp}_T(G).
\]
In particular, if $G$ is $(T,\rho,\sigma)$-visibility-rich and
\[
 m \ge \frac{\sigma}{\rho}\ln\frac{1}{\varepsilon},
\]
then
\[
\mathsf{Occ}_{m,T}(G) \ge (1-\varepsilon)\sigma.
\]
\end{corollary}

\begin{corollary}[Visible coverage time]
For $0 < \alpha < 1$, define the $\alpha$-coverage time $\tau_{\alpha,T}(G)$ as the least $m$ such that
\[
\mathsf{Occ}_{m,T}(G) \ge \alpha\mathsf{VSupp}_T(G).
\]
Then
\[
\tau_{\alpha,T}(G) \le \frac{\mathsf{VSupp}_T(G)}{\mathsf{Vis}_T(G)}\ln\frac{1}{1-\alpha}.
\]
In particular, the visible half-life $\tau_{1/2,T}(G)$ is at most
\[
\frac{\mathsf{VSupp}_T(G)}{\mathsf{Vis}_T(G)}\ln 2.
\]
\end{corollary}

\begin{remark}
The ratio $\mathsf{VSupp}_T(G)/\mathsf{Vis}_T(G)$ is the characteristic visible timescale: it controls how quickly the occupancy curve approaches saturation and gives a direct diagnostic for generators whose valid outputs are technically present but practically hard to discover.
\end{remark}

\begin{remark}
One can also view the observed occupancy $\widehat{\mathsf{Occ}}_{m,T}(G)$ as a 1-Lipschitz function of the sample sequence. In fact, it is also 1-self-bounding, so the visible occupancy profile admits both bounded-differences concentration and variance control~\cite{McDiarmid1989,Warnke2016,BoucheronLugosiMassart2009,McDiarmidReed2006}. This suggests a next-step refinement: high-probability versions of the coverage and discovery statements rather than only expectations.
\end{remark}

\begin{proposition}[Visible occupancy concentration]
Let $\widehat{\mathsf{Occ}}_{m,T}(G)$ be the random number of distinct outputs in $\mathsf{Heavy}_T(G)$ observed among $X_1,\dots,X_m$, where $X_1,\dots,X_m$ are independent samples from $D_G$. Then
\[
\mathbb{E}[\widehat{\mathsf{Occ}}_{m,T}(G)] = \mathsf{Occ}_{m,T}(G).
\]
Moreover, for every $t > 0$,
\[
\Pr\bigl[|\widehat{\mathsf{Occ}}_{m,T}(G) - \mathsf{Occ}_{m,T}(G)| \ge t\bigr] \le 2\exp\Bigl(-\frac{2t^2}{m}\Bigr).
\]
In particular, with probability at least $1-\delta$,
\[
\widehat{\mathsf{Occ}}_{m,T}(G) \ge \mathsf{Occ}_{m,T}(G) - \sqrt{\frac{m}{2}\ln\frac{2}{\delta}}.
\]
\end{proposition}

\begin{proposition}[Visible occupancy self-bounding]
Let $\widehat{\mathsf{Occ}}_{m,T}(G)$ be as above. Then it is a 1-self-bounding function of the sample sequence. In particular,
\[
\mathrm{Var}[\widehat{\mathsf{Occ}}_{m,T}(G)] \le \mathbb{E}[\widehat{\mathsf{Occ}}_{m,T}(G)] = \mathsf{Occ}_{m,T}(G).
\]
Consequently, the fluctuations of visible occupancy are Poisson-scale in the sparse regime.
\end{proposition}

\bibliographystyle{plain}
\bibliography{references}

\end{document}
